\section{Matrix Equations}

\begin{frame}{Warm-up: The Row-Reduction Algorithm}
Find the inverse of the coefficient matrix from the next example:
\[
A=\begin{bmatrix}-7&3\\5&-2\end{bmatrix}.
\]
\pause
\small
\[
\left[\begin{array}{cc|cc}-7&3&1&0\\5&-2&0&1\end{array}\right]
\xrightarrow{R_1\leftarrow R_1+R_2}
\left[\begin{array}{cc|cc}-2&1&1&1\\5&-2&0&1\end{array}\right]
\]
\[
\xrightarrow{R_2\leftarrow R_2+2R_1}
\left[\begin{array}{cc|cc}-2&1&1&1\\1&0&2&3\end{array}\right]
\xrightarrow{R_1\leftrightarrow R_2}
\left[\begin{array}{cc|cc}1&0&2&3\\-2&1&1&1\end{array}\right]
\]
\[
\xrightarrow{R_2\leftarrow R_2+2R_1}
\left[\begin{array}{cc|cc}1&0&2&3\\0&1&5&7\end{array}\right].
\]
\pause
Therefore
\[
A^{-1}=\begin{bmatrix}2&3\\5&7\end{bmatrix}.
\]
\end{frame}

\begin{frame}{Solving a $2\times2$ System with $A^{-1}$}
Solve
\[
-7x_1+3x_2=2,\qquad 5x_1-2x_2=1.
\]
\pause
In matrix form,
\[
A\vect x=\vect b,\qquad
A=\begin{bmatrix}-7&3\\5&-2\end{bmatrix},\quad
\vect b=\begin{bmatrix}2\\1\end{bmatrix}.
\]
\pause
Multiplying on the left by the inverse gives
\[
\vect x=A^{-1}\vect b
=\begin{bmatrix}2&3\\5&7\end{bmatrix}
\begin{bmatrix}2\\1\end{bmatrix}
=\begin{bmatrix}7\\17\end{bmatrix}.
\]
\end{frame}

\begin{frame}{Why the Solution Is Unique}
Suppose $A$ is invertible and $A\vect x=\vect b$. Then
\[
A^{-1}A\vect x=A^{-1}\vect b
\quad\Longrightarrow\quad
\vect x=A^{-1}\vect b.
\]
\pause
If $\vect w$ is another solution, then
\[
A\vect w=\vect b
\quad\Longrightarrow\quad
\vect w=A^{-1}\vect b.
\]
Thus $\vect w=\vect x$. Any solution must equal $A^{-1}\vect b$.
\end{frame}

\begin{frame}{Unique Solutions from Invertibility}
\theorembox{Theorem 5}{If $A$ is an invertible $n\times n$ matrix, then for every $\vect b\in\R^n$, the equation $A\vect x=\vect b$ has the unique solution
\[
\vect x=A^{-1}\vect b.
\]}
\medskip
\begin{block}{Homogeneous special case}
If $A$ is invertible, then $A\vect x=\vect0$ has only the trivial solution $\vect x=\vect0$.
\end{block}
\end{frame}

\begin{frame}{Solution Types of Linear Systems}
For a general $m\times n$ matrix $A$, the system $A\vect x=\vect b$ may have
\[
\text{no solution},\qquad \text{one solution},\qquad\text{or infinitely many solutions}.
\]
\pause
If $A$ is square and invertible, every $A\vect x=\vect b$ has exactly one solution:
\[
\vect x=A^{-1}\vect b.
\]
\pause
The homogeneous system $A\vect x=\vect0$ always has the trivial solution. It has either only that solution or infinitely many solutions.
\end{frame}

\begin{frame}{Example: The Inverse-Matrix Algorithm}
Solve
\[
\begin{cases}
x_1+2x_2-x_3=-5,\\
2x_1-x_2+x_3=6,\\
x_1-x_2-3x_3=-3.
\end{cases}
\]
\pause
We will write the system as $A\vect x=\vect b$ and compute $\vect x=A^{-1}\vect b$.
\end{frame}

\begin{frame}{Example: Matrix Form}
\[
A=\begin{bmatrix}
1&2&-1\\2&-1&1\\1&-1&-3
\end{bmatrix},\qquad
\vect x=\begin{bmatrix}x_1\\x_2\\x_3\end{bmatrix},\qquad
\vect b=\begin{bmatrix}-5\\6\\-3\end{bmatrix}.
\]
The inverse, found by row reduction, is
\[
A^{-1}=\frac1{19}
\begin{bmatrix}
4&7&1\\7&-2&-3\\-1&3&-5
\end{bmatrix}.
\]
\end{frame}

\begin{frame}{Example: Solution}
\[
\vect x=A^{-1}\vect b
=\frac1{19}
\begin{bmatrix}
4&7&1\\7&-2&-3\\-1&3&-5
\end{bmatrix}
\begin{bmatrix}-5\\6\\-3\end{bmatrix}
\]
\pause
\[
=\frac1{19}\begin{bmatrix}19\\-38\\38\end{bmatrix}
=\begin{bmatrix}1\\-2\\2\end{bmatrix}.
\]
Thus $x_1=1$, $x_2=-2$, and $x_3=2$.
\end{frame}

\begin{frame}{General Matrix Equations $AX=B$}
Suppose $A$ is $n\times n$ and $X,B$ are $n\times s$. If $A$ is invertible,
\[
A^{-1}AX=A^{-1}B
\quad\Longrightarrow\quad
\boxed{X=A^{-1}B}.
\]
\begin{block}{Conditions for the inverse-matrix algorithm}
\begin{enumerate}
  \item The coefficient matrix $A$ is square.
  \item The coefficient matrix $A$ is invertible.
\end{enumerate}
\end{block}
\end{frame}

\begin{frame}[shrink=8]{Three Common Matrix Equations}
% \begin{columns}[T]
% \column{.28\textwidth}\centering
\begin{itemize}
  \item $AX=B$ \pause $\sim X=A^{-1}B$\pause
  \item $XA=B$ \pause $\sim X=BA^{-1}$\pause
  \item $AXB=C$ \pause $\sim X=A^{-1}CB^{-1}$
\end{itemize}
% \[AX=B,\qquad X=A^{-1}B\]\\
% \column{.28\textwidth}\centering
% \[XA=B,\qquad X=BA^{-1}\]\\
% \column{.40\textwidth}\centering
% \[AXB=C,\qquad X=A^{-1}CB^{-1}\]\\
% \end{columns}
% \pause
% For example, solve $AX=B$ with
% \[
% A=\begin{bmatrix}0&-1&1\\1&0&0\\1&0&2\end{bmatrix},\qquad
% B=\begin{bmatrix}0&-1\\2&-5\\0&1\end{bmatrix}.
% \]
% Since
% \[
% A^{-1}=\begin{bmatrix}0&1&0\\-1&-\frac12&\frac12\\0&-\frac12&\frac12\end{bmatrix},
% \quad
% X=A^{-1}B=\begin{bmatrix}2&-5\\-1&4\\-1&3\end{bmatrix}.
% \]
\end{frame}

\begin{frame}[shrink=8]{Exercises: Matrix Equations I}
\textbf{1.} Solve $XA=B$ for
\[
A=\begin{bmatrix}1&0&-2\\2&4&-1\\0&1&1\end{bmatrix},\qquad
B=\begin{bmatrix}1&0&2\\2&4&3\end{bmatrix}.
\]
\pause
\[
A^{-1}=\begin{bmatrix}5&-2&8\\-2&1&-3\\2&-1&4\end{bmatrix},
\quad X=BA^{-1}=\begin{bmatrix}9&-4&16\\8&-3&16\end{bmatrix}.
\]
\pause
\textbf{2.} Solve
\[
\begin{bmatrix}1&4\\0&1\end{bmatrix}X
\begin{bmatrix}1&0\\3&1\end{bmatrix}
=\begin{bmatrix}-1&1\\1&0\end{bmatrix}.
\]
Multiplying by the two inverses gives $X=\begin{bmatrix}-8&1\\1&0\end{bmatrix}$.
\end{frame}

\begin{frame}{Exercise: Collect the $X$ Terms}
\textbf{3.} Solve $AX=A+2X$ for
\[
A=\begin{bmatrix}4&0&0\\0&1&0\\0&0&1\end{bmatrix}.
\]
\pause
\[(A-2I)X=A,\qquad X=(A-2I)^{-1}A.\]
\end{frame}

\begin{frame}{Exercise: Multiply on the Correct Side}
\textbf{4.} Solve $A^2-XA-I=0$, where
\[
A=\begin{bmatrix}1&0&3\\2&1&1\\-1&0&1\end{bmatrix}.
\]
\pause
\[
X=(A^2-I)A^{-1}=A-A^{-1}
=\begin{bmatrix}.75&0&3.75\\2.75&0&-.25\\-1.25&0&.75\end{bmatrix}.
\]
\end{frame}

\begin{frame}{Exercise: Isolating an Unknown Matrix}
Suppose $A,B,C$ are invertible $n\times n$ matrices and $A=B(D-I_n)C$. Solve $D$.
\pause
Multiplying on the left by $B^{-1}$ and on the right by $C^{-1}$ gives
\[
B^{-1}AC^{-1}=D-I_n.
\]
Therefore
\[
\boxed{D=B^{-1}AC^{-1}+I_n}.
\]
The order of the factors must be preserved.
\end{frame}

\begin{frame}[shrink=7]{Exercise: A Matrix Equation with a Transpose}
Let
\[
ABA^T=2BA^T+I,\qquad
A=\begin{bmatrix}1&0&0\\0&1&2\\0&0&1\end{bmatrix}.
\]. Solve B.
\pause
Since $2BA^T=2IBA^T$,
\[
(A-2I)BA^T=I,
\quad B=(A-2I)^{-1}(A^T)^{-1}=\bigl[A^T(A-2I)\bigr]^{-1}.
\]
\pause
\[
A^T(A-2I)=\begin{bmatrix}-1&0&0\\0&-1&2\\0&-2&3\end{bmatrix},
\quad
\boxed{B=\begin{bmatrix}-1&0&0\\0&3&-2\\0&2&-1\end{bmatrix}}.
\]
\end{frame}
